\section{Representation Spaces}

A fundamental aspect of machine learning systems is how information is represented internally.
Since real-world information such as text, images, and signals is not naturally expressed in a format suitable for mathematical operations, it must be transformed into a structured representation.

A representation space defines how information is encoded and how relationships between individual elements are expressed. Depending on the chosen representation, information can be organized as discrete states or continuous vectors. These approaches differ in how information is stored and how relationships between individual representations can be modeled.

In general, a representation function maps an input sample $x$ from an input space $\mathcal{X}$ into a representation space $\mathcal{Z}$:

\begin{equation}
    f: \mathcal{X} \rightarrow \mathcal{Z}
\end{equation}

where $\mathcal{Z}$ defines the space in which information is stored and compared.

\section{Representation Learning}

While representation spaces define how information can be encoded and compared, the quality of a representation strongly depends on how well it captures relevant characteristics of the underlying data. Representation learning refers to the process of automatically learning such representations from data rather than relying on manually designed features.

Given a dataset of input samples, a representation learning model learns a mapping function that transforms raw inputs into representations where relevant relationships are preserved and irrelevant variations are minimized. The objective is to learn a representation space in which samples with similar semantic meaning are located closer together, while dissimilar samples are separated.

In the context of natural language processing, representation learning enables the transformation of textual information into dense vector representations, commonly referred to as embeddings. Modern Transformer-based models learn contextualized representations in which the meaning of a token or sequence depends on its surrounding context. These representations can subsequently be adapted through fine-tuning on domain-specific data to better capture task-specific relationships.

For entity linking tasks, representation learning is particularly important because the model must learn semantic relationships between textual mentions and their corresponding entities. Rather than relying on exact lexical overlap, learned representations enable the identification of semantically equivalent references that may use different surface forms.


\subsection{Discrete Representation Spaces}

In discrete representation spaces, information is represented through a finite set of distinct states:

\begin{equation}
    \mathcal{Z}_{D} = \{z_1, z_2, ..., z_n\}
\end{equation}

where each representation $z_i$ corresponds to a specific symbol, category, or pattern. Relationships between different states are not inherently represented, meaning that two states are either identical or distinct.

This type of representation is commonly used in symbolic systems and early associative memory models. For example, classical Hopfield networks represent memories as discrete binary patterns:

\begin{equation}
    \xi_i \in \{-1,+1\}^{N}
\end{equation}

where each stored pattern corresponds to a stable state of the network~\cite{hopfield1982}.

While discrete representations allow for clear separation between individual states, they provide limited capability for expressing gradual relationships or similarities between different representations.



\subsection{Continuous Representation Spaces}

Continuous representation spaces represent information as points within a high-dimensional vector space:

\begin{equation}
    \mathcal{Z}_{C} \subseteq \mathbb{R}^{d}
\end{equation}

where each representation is encoded as a vector:

\begin{equation}
    z = f(x) = [z_1,z_2\ldots,z_d]^{T}
\end{equation}

Instead of assigning isolated states to individual concepts or patterns, information is encoded through numerical vectors, where relationships between representations can be expressed through geometric properties such as distance or similarity.

The use of distributed continuous representations has become a central concept in machine learning, where information is represented through patterns of numerical values rather than isolated symbolic states~\cite{hinton1986learning}. Similar information can be represented by nearby points within the vector space, allowing the representation space itself to capture relationships between stored information.

Within associative memory systems, continuous representation spaces enable memories to be represented as high-dimensional vectors rather than discrete states. A memory set can therefore be expressed as:

\begin{equation}
    M = \{m_1,m_2\ldots,m_k\}, \quad m_i \in \mathbb{R}^{d}
\end{equation}

where each $m_i$ represents an individual stored memory within the continuous representation space. This provides a more flexible representation of information and forms the basis for modern approaches to associative memory and machine learning systems~\cite{ramsauer2021hopfield}.

In this work, continuous representation spaces serve as the foundation for storing and accessing memories.